\documentclass[10pt, letterpaper]{article}

% Packages:
\usepackage[
    ignoreheadfoot, % set margins without considering header and footer
    top=2 cm, % seperation between body and page edge from the top
    bottom=2 cm, % seperation between body and page edge from the bottom
    left=2 cm, % seperation between body and page edge from the left
    right=2 cm, % seperation between body and page edge from the right
    footskip=1.0 cm, % seperation between body and footer
    % showframe % for debugging 
]{geometry} % for adjusting page geometry
\usepackage{titlesec} % for customizing section titles
\usepackage{tabularx} % for making tables with fixed width columns
\usepackage{array} % tabularx requires this
\usepackage[dvipsnames]{xcolor} % for coloring text
\definecolor{primaryColor}{RGB}{0, 0, 0} % define primary color
\usepackage{enumitem} % for customizing lists
\usepackage{fontawesome5} % for using icons
\usepackage{amsmath} % for math
\usepackage[
    pdftitle={John Doe's CV},
    pdfauthor={John Doe},
    pdfcreator={LaTeX with RenderCV},
    colorlinks=true,
    urlcolor=primaryColor
]{hyperref} % for links, metadata and bookmarks
\usepackage[pscoord]{eso-pic} % for floating text on the page
\usepackage{calc} % for calculating lengths
\usepackage{bookmark} % for bookmarks
\usepackage{lastpage} % for getting the total number of pages
\usepackage{changepage} % for one column entries (adjustwidth environment)
\usepackage{paracol} % for two and three column entries
\usepackage{ifthen} % for conditional statements
\usepackage{needspace} % for avoiding page brake right after the section title
\usepackage{iftex} % check if engine is pdflatex, xetex or luatex

% Ensure that generate pdf is machine readable/ATS parsable:
\ifPDFTeX
    \input{glyphtounicode}
    \pdfgentounicode=1
    \usepackage[T1]{fontenc}
    \usepackage[utf8]{inputenc}
    \usepackage{lmodern}
\fi

\usepackage{charter}

% Some settings:
\raggedright
\AtBeginEnvironment{adjustwidth}{\partopsep0pt} % remove space before adjustwidth environment
\pagestyle{empty} % no header or footer
\setcounter{secnumdepth}{0} % no section numbering
\setlength{\parindent}{0pt} % no indentation
\setlength{\topskip}{0pt} % no top skip
\setlength{\columnsep}{0.15cm} % set column seperation
\pagenumbering{gobble} % no page numbering

\titleformat{\section}{\needspace{4\baselineskip}\bfseries\large}{}{0pt}{}[\vspace{1pt}\titlerule]

\titlespacing{\section}{
    % left space:
    -1pt
}{
    % top space:
    0.3 cm
}{
    % bottom space:
    0.2 cm
} % section title spacing

\renewcommand\labelitemi{$\vcenter{\hbox{\small$\bullet$}}$} % custom bullet points
\newenvironment{highlights}{
    \begin{itemize}[
        topsep=0.10 cm,
        parsep=0.10 cm,
        partopsep=0pt,
        itemsep=0pt,
        leftmargin=0 cm + 10pt
    ]
}{
    \end{itemize}
} % new environment for highlights


\newenvironment{highlightsforbulletentries}{
    \begin{itemize}[
        topsep=0.10 cm,
        parsep=0.10 cm,
        partopsep=0pt,
        itemsep=0pt,
        leftmargin=10pt
    ]
}{
    \end{itemize}
} % new environment for highlights for bullet entries

\newenvironment{onecolentry}{
    \begin{adjustwidth}{
        0 cm + 0.00001 cm
    }{
        0 cm + 0.00001 cm
    }
}{
    \end{adjustwidth}
} % new environment for one column entries

\newenvironment{twocolentry}[2][]{
    \onecolentry
    \def\secondColumn{#2}
    \setcolumnwidth{\fill, 4.5 cm}
    \begin{paracol}{2}
}{
    \switchcolumn \raggedleft \secondColumn
    \end{paracol}
    \endonecolentry
} % new environment for two column entries

\newenvironment{threecolentry}[3][]{
    \onecolentry
    \def\thirdColumn{#3}
    \setcolumnwidth{, \fill, 4.5 cm}
    \begin{paracol}{3}
    {\raggedright #2} \switchcolumn
}{
    \switchcolumn \raggedleft \thirdColumn
    \end{paracol}
    \endonecolentry
} % new environment for three column entries

\newenvironment{header}{
    \setlength{\topsep}{0pt}\par\kern\topsep\centering\linespread{1.5}
}{
    \par\kern\topsep
} % new environment for the header

\newcommand{\placelastupdatedtext}{% \placetextbox{<horizontal pos>}{<vertical pos>}{<stuff>}
  \AddToShipoutPictureFG*{% Add <stuff> to current page foreground
    \put(
        \LenToUnit{\paperwidth-2 cm-0 cm+0.05cm},
        \LenToUnit{\paperheight-1.0 cm}
    ){\vtop{{\null}\makebox[0pt][c]{
        \small\color{gray}\textit{Last updated in September 2024}\hspace{\widthof{Last updated in September 2024}}
    }}}%
  }%
}%

% save the original href command in a new command:
\let\hrefWithoutArrow\href

% new command for external links:


\begin{document}
    \newcommand{\AND}{\unskip
        \cleaders\copy\ANDbox\hskip\wd\ANDbox
        \ignorespaces
    }
    \newsavebox\ANDbox
    \sbox\ANDbox{$|$}

    \begin{header}
        \fontsize{25 pt}{25 pt}\selectfont Ariq Fadhilah Jeha

        \vspace{5 pt}
        \normalsize
        \mbox{Depok}%
        \kern 5.0 pt%
        \AND%
        \kern 5.0 pt%
        \mbox{\hrefWithoutArrow{mailto:ariqfadhh@gmail.com}{ariqfadhh@gmail.com}}%
        \kern 5.0 pt%
        \AND%
        \kern 5.0 pt%
        \mbox{\hrefWithoutArrow{tel:+62895622994052}{+62-8956-2299-4052}}%
        \kern 5.0 pt%
        \AND%
        \kern 5.0 pt%
        \mbox{\hrefWithoutArrow{https://linkedin.com/in/ariqfadhilah}{linkedin.com/in/ariqfadhilah}}%
        \kern 5.0 pt%
        \AND%
        \kern 5.0 pt%
        \mbox{\hrefWithoutArrow{https://github.com/Ariqfadh}{github.com/Ariqfadh}}%
    \end{header}

    \vspace{5 pt - 0.1 cm}

    \section{Background}
    
    \begin{onecolentry}
        I am a Biomedical Engineering student with a strong interest in Integrated Circuit (IC) design, spanning both analog and digital domains. I have hands-on experience in both analog and digital , from functional verification untill post-layout simulation using Cadence.
    \end{onecolentry}

    \vspace{0.1cm}

    \section{Education}
        
        \begin{twocolentry}{
            August 2022 – Present 
        }
            \textbf{Gadjah Mada University}, Biomedical Engineering\end{twocolentry}

        \vspace{0.10 cm}
        \begin{onecolentry}
            \begin{highlights}
                \item GPA: 3.39/4.0\\
                \item \textbf{Coursework:} Fundamental Electronics, Biomedical Electronics, Digital Engineering, Microcontroller System , VLSI Development, Digital System Development.
            \end{highlights}
        \end{onecolentry}

    
\section{Experience}

    %===================Kerja Praktek
    \begin{twocolentry}{
        Jan 2025 – Feb 2025
    }
        \textbf{Research Intern}, Beehive Drones -- Jogjakarta
    \end{twocolentry}
    
    \vspace{0.1 cm}
    
    \begin{onecolentry}
        \begin{highlights}
            \item Developed a FAST keypoint detection module using Vitis HLS.  
            \item Designed and verified the accelerator at C/C++ and HLS simulation level.  
        \end{highlights}
    \end{onecolentry}
    
    \vspace{0.2 cm}
    
    %===================Riset intern
    \begin{twocolentry}{
        May 2025 – May 2026
    }
        \textbf{Research Assistant}, DTETI -- Jogjakarta
    \end{twocolentry}
    
    \vspace{0.10 cm}
    
    \begin{onecolentry}
        \begin{highlights}
            \item Developing custom AXI4-Lite peripherals for a RISC-V-based microcontroller, including SPI and UART interfaces.
            \item Designed and verified Verilog RTL using simulation and custom testbenches.
        \end{highlights}
    \end{onecolentry}
        
        
    \section{Projects}
            
    \begin{twocolentry}{}
        \textbf{1. Phase-Locked Loop In MICS Band}
    \end{twocolentry}
    
    \vspace{0.10 cm}
    \begin{onecolentry}
        \begin{highlights}
            \item Designed and implemented a full-custom Integer-N PLL in CMOS GPDK045 for 402--405 MHz MICS-band applications, including PFD, charge pump, loop filter, current-starved VCO, and frequency divider.  
            \item Completed full-custom layout, DRC/LVS verification, parasitic extraction, and post-layout simulation, achieving 402--405 MHz output with 15--39 $\mu$s settling time and 214.85--230.2 $\mu$W power consumption.
        \end{highlights}
    \end{onecolentry}
    
    \vspace{0.3cm}

    \begin{twocolentry}{}
        \textbf{2. Ambulatory ECG System }
    \end{twocolentry}
    
    \vspace{0.10 cm}
    \begin{onecolentry}
        \begin{highlights}
            \item Developed an ambulatory ECG acquisition system using STM32F4 (300 Hz sampling rate).  
            \item Implemented timer-driven UART acquisition with interrupt + DMA and circular buffering.  
            \item Transmitted real-time ECG data via Bluetooth.  
        \end{highlights}
    \end{onecolentry}
        
    \vspace{0.3cm}
    
    \begin{twocolentry}{}
        \textbf{3. AXI4-Lite SPI \& UART Peripheral Design}
    \end{twocolentry}
    
    \vspace{0.30cm}
    
    \begin{onecolentry}
        \begin{highlights}
            \item Designed SPI master and UART communication peripherals.  
            \item Implemented AXI4-Lite wrappers with memory-mapped control and status registers.  
            \item Integrated FIFO buffering for UART data handling.  
            \item Verified functionality through RTL simulation and testbench-based validation.  
        \end{highlights}
    \end{onecolentry}
    
        
\section{Technologies}

\begin{onecolentry}
    \textbf{Languages:} Tcl, Verilog, SystemVerilog, Python, C++
\end{onecolentry}

\vspace{0.1 cm}

\begin{onecolentry}
    \textbf{EDA \& Tools:} Cadence Virtuoso, Spectre ADE, Genus, Innovus, Tempus, Xilinx Vivado, Vitis
\end{onecolentry}

\vspace{0.1 cm}

\begin{onecolentry}
    \textbf{Environment:} Linux (CLI, shell scripting)
\end{onecolentry}
\end{document}